\documentclass[12pt]{article}
\usepackage[T1]{fontenc}
\usepackage[utf8]{inputenc}
\usepackage{lmodern}
\usepackage{textcomp}
\usepackage{enumitem}
\usepackage{geometry}
\usepackage{graphicx}
\usepackage{amsmath, amssymb}
\geometry{margin=1in}
\begin{document}
\begin{center}
Economics - Undergraduate Level Assessment 7 \
Total Marks: 40 \
Answer all questions.
\end{center}

\section*{Multiple Choice (10 marks)}
\begin{enumerate}[label=\arabic*.]
\item What is a monopoly?
\begin{enumerate}[label=(\alph*)]
    \item A monopoly is a company that has no control over the market.
    \item A monopoly is a government-run entity that controls all aspects of a market.
    \item A monopoly is a business that can set its own prices.
    \item A monopoly is a firm which is the sole producer of a good.
    \item A monopoly is a market structure with only two firms competing against each other.
\end{enumerate}
\item What is the primary objective of management in managerial economics?
\begin{enumerate}[label=(\alph*)]
    \item maximization of the firm's value
    \item minimization of costs
    \item maximization of market share
    \item minimization of environmental impact
    \item maximization of social welfare
\end{enumerate}
\item The price elasticity of demand for a product is greater if
\begin{enumerate}[label=(\alph*)]
    \item the proportion of the good of the consumer's budget is high.
    \item the number of substitute products is limited.
    \item the product is unique and has no substitutes.
    \item the consumer has a high income.
    \item the period of time to respond to a price change is short.
\end{enumerate}
\item At the peak of a typical business cycle which of the following is likely the greatest threat to the macroeconomy?
\begin{enumerate}[label=(\alph*)]
    \item Unemployment
    \item Overproduction of goods and services
    \item Bankruptcy
    \item Reduced consumer spending
    \item Decreased interest rates
\end{enumerate}
\item If people expect the price of a particular product to increase in the near future
\begin{enumerate}[label=(\alph*)]
    \item this will increase the supply of the product.
    \item this will cause the producer to decrease the supply of the product.
    \item this will not affect the demand for the product now or later.
    \item this will decrease the supply of the product.
    \item this will increase the demand for the product.
\end{enumerate}
\end{enumerate}

\section*{True / False (10 marks)}
\textit{Answer True or False}
\begin{enumerate}[label=\arabic*.]
\setcounter{enumi}{5}
\item True or False: In the presence of a positive externality, the marginal private cost exceeds the marginal social benefit at the market quantity.
\item True or False: The final model using the general to specific approach may lack theoretical interpretation.
\item True or False: The GDP Deflator is less reliable than the Consumer Price Index (CPI) in measuring overall economic inflation.
\item True or False: An increase in the standard of living indicates that output has increased proportionally more than the population.
\item True or False: The removal of a protective tariff on imported steel improves allocative efficiency in the market.
\end{enumerate}

\section*{Long Form Response (20 marks)}
\textit{Answer each question with a clear and complete explanation.}
\begin{enumerate}[label=\arabic*.]
\setcounter{enumi}{10}
\item Discuss the differences in price determination between centrally planned and capitalist economies, and analyze the implications of these differences on the economic functions of prices in each system.
\item Discuss the key characteristics of British socialism, focusing on its evolutionary approach, government involvement in industries, and social welfare initiatives. Analyze their economic implications.
\end{enumerate}

\vfill
\noindent\textit{End of Paper}
\end{document}